A Neural ODE \cite{chen2018neuralode} learns an unknown dynamical system
$\dot{\mathbf{u}}(t) = \mathbf{F}(\mathbf{u}(t))$ by fitting a vector field
$\mathbf{f}_\theta(\mathbf{u}(t))$ and integrating it numerically. The usual
choice for $\mathbf{f}_\theta$ is a multilayer perceptron (MLP), with fixed
activations on the nodes and all learnable structure in the linear weights,
$\mathbf{y} = \sigma(\mathbf{W}\mathbf{x}+\mathbf{b})$. A Kolmogorov-Arnold
Network (KAN) \cite{liu2024kan} flips this around: every \emph{edge} carries
its own learnable univariate function $\phi(x)$, and each node just sums its
incoming edges, $y_j = \sum_i \phi_{i,j}(x_i)$. The idea comes from the
Kolmogorov-Arnold representation theorem, which says that any continuous
multivariate function on a bounded domain can be written as a finite
composition of continuous univariate functions and addition,
\begin{equation}
f(\mathbf{x}) = \sum_{q=1}^{2n+1}\Phi_q\!\left(\sum_{p=1}^{n}\phi_{q,p}(x_p)\right),
\label{eq:kolmogorov-arnold}
\end{equation}
so no nonlinearity has to be fixed in advance; the network learns the shape
of each univariate function. In practice each edge function is a fixed
smooth residual plus a learnable expansion in a chosen basis. In the form we
implement,
\begin{equation}
\phi(x) = w_b\, b(x) + \sum_{i=1}^{G} c_i\, B_i\big(\mathrm{norm}(x)\big),
\label{eq:kan-edge}
\end{equation}
with $b(x)=\mathrm{silu}(x) = x/(1+e^{-x})$ and $\mathrm{norm}(x)=\tanh x$.
Here $B_1,\dots,B_G$ are $G$ basis functions attached to a uniform grid of
knots $z_g\in[-1,1]$ with spacing $h=2/(G-1)$, the $c_i$ are trainable
coefficients, and $w_b$ weights the residual. (The original \texttt{pykan}
form carries $G+k$ spline functions and a separate scale $w_s$ on the
expansion; with one coefficient per basis function that scale is redundant,
so we drop it.) The choice of $B_i$ is a design decision rather than a
mathematical necessity. A Gaussian radial basis function (RBF) grid,
$B_i(u)=\exp(-(u-z_i)^2/2h^2)$, gives each edge smooth support that decays
quickly but never reaches zero. A degree-$k$ B-spline built by the
Cox-de Boor recursion instead gives each edge compact local support, exactly
zero outside its own knot span.

\subsection{The base paper: KAN-ODEs}
\label{sec:base-paper}

Koenig et al.\ \cite{koenig2024kanode} put this construction inside a Neural
ODE's vector field, $\dot{\mathbf{u}}=\mathbf{f}_\theta(\mathbf{u})$ with
every edge of $\mathbf{f}_\theta$ built from Equation~\ref{eq:kan-edge}, and
obtain the trajectory by numerically integrating that field forward in time.
On the Lotka-Volterra predator-prey system they report two central claims:
a KAN-ODE with 240 parameters extrapolates the system's oscillation
accurately from a short training window, and it reaches a target training
loss in roughly $10^4$ epochs where a parameter-matched (252-parameter)
MLP-based Neural ODE needs $10^5$, a tenfold training-speed advantage. The
paper also applies KAN-ODEs to other systems, including inferring hidden
physics; we focus on its Lotka-Volterra benchmark.

To get these results the paper commits to three design decisions, which we
treat as free variables:
\begin{enumerate}
\item the \textbf{integrator}: an adaptive fifth-order Runge-Kutta method,
  Tsit5 \cite{tsitouras2011tsit5};
\item the \textbf{basis}: Gaussian RBF;
\item the \textbf{differentiation method} used to get gradients through the
  integrator: continuous adjoint sensitivity, chosen for its constant memory
  cost in trajectory length.
\end{enumerate}
The basis choice is not something we invented to vary. The original KAN
paper \cite{liu2024kan} used cubic B-splines by default, and the base paper
switched to Gaussian RBF for the ODE setting because its derivative never
vanishes, which matters once the edge function sits inside a differentiated
integrator rather than a static regression. We restore B-spline and add five
further bases, so the trade-off behind that switch is measured rather than
assumed.

\subsection{Why training needs to be studied, not just the final error}

Besides Lotka-Volterra, we use two textbook systems, the SIR epidemic model
and a damped pendulum, chosen because their governing equations are simple
and well understood. Any training failure can then be blamed on the
architecture or the optimization rather than on ambiguous target dynamics.
The integrator's interaction with training is not predictable from classical
numerical analysis alone, because the integrator sits inside a
backpropagation graph rather than a standalone forward solve: a solver that
is perfectly stable in the classical sense can still train badly, train to a
wrong answer, or fail to generalize. Section~\ref{sec:numerics} makes this
concrete. A network trained through Forward Euler learns the \emph{inverse
modified equation} of backward error analysis rather than the physical
field, and every trained model, KAN or MLP, replaces Lotka-Volterra's family
of neutral orbits with a single attracting limit cycle. Diagnosing behaviour
like this means treating training itself as an object of study.

\subsection{Summary of our work}
\label{sec:work-summary}

Table~\ref{tab:work-summary} lists every study we ran, why we ran it, what we
found, and where it is reported. In short, our objectives were to:
(i)~reproduce the base paper's Lotka-Volterra benchmark with every
integrator and basis implemented and verified from its defining formula;
(ii)~vary each of the paper's three fixed design choices and measure what
it actually costs or buys; (iii)~check whether the recipe transfers to other
dynamical systems, and fix it where it does not; and (iv)~test whether our
conclusions survive a five times larger training budget.

\begin{table*}[t]
\centering
\footnotesize
\renewcommand{\arraystretch}{1.18}
\caption{Summary of the studies in this report.}
\label{tab:work-summary}
\begin{tabularx}{\textwidth}{@{}L{2.7cm} L{4.6cm} Y C{0.7cm}@{}}
\toprule
\hdrrow \thh{Study} & \thh{Why we ran it} & \thh{What we observed} & \thh{Sec.}\\
\midrule
Baseline reproduction: KAN-ODE vs.\ MLP-ODE & Test the base paper's two headline claims (accuracy and training speed) at matched parameter count. & KAN-ODE extrapolates $5.4\times$ better at $10^4$ epochs, a gap explained by a $48\times$ smaller period error. The MLP trains $4.1\times$ \emph{faster}. The paper's literal tanh MLP trains only under the paper's own learning rate and initialization. & \ref{sec:kan-vs-mlp}\\
Integrator ablation & The paper fixes Tsit5; does solver order matter for training? & All six solvers pass their order conditions to round-off. Above order two, accuracy is flat: Midpoint matches Tsit5 within $0.3\%$ at $29\%$ of the cost. The Euler model learns the inverse modified equation, and no model transfers across solver families. & \ref{sec:solver-results}\\
Basis-function ablation & The paper switched from B-spline to RBF; is that justified? & B-spline and RBF are tied, with RBF $3.9\times$ cheaper. Global polynomials extrapolate $4$--$5\times$ worse. Locality, not conditioning, separates the bases. & \ref{sec:basis-results}\\
Step-size and noise sensitivity & How sensitive is the fit to sampling density and observation noise? & Step size acts as a data-density choice (non-monotone). Extrapolation MSE scales as $12.9\,\sigma^{2.02}$ with noise. & \ref{sec:dt-noise}\\
What the trained fields learned & Does low trajectory error mean the right vector field? & No. Every model is an attracting limit cycle around a spurious unstable equilibrium, accurate only in a thin tube around the training orbit. & \ref{sec:numerics}\\
Cross-domain failures and fixes & Does the Lotka-Volterra recipe transfer to the pendulum and SIR? & A longer training window fixes the pendulum. SIR's $19\times$ scale mismatch at initialization is fixed by time rescaling and a conservation projection; adding a vanishing-field gate gives $R^2$: $-20.97\to+0.998$. The KAN advantage widens to $40.6\times$ at twice the horizon. & \ref{sec:crossdomain}\\
Adjoint vs.\ direct autograd & The paper uses the adjoint for its memory saving; does that matter at our scale? & The adjoint uses less memory at every length tested but is $1.9$--$2.1\times$ slower, so we use direct autograd. & \ref{sec:track-a}\\
Gradient-norm dynamics & Is Euler's convergence floor caused by noisier gradients? & No ($\rho=-0.088$, $p=0.87$); gradient spike bursts are shared across all solvers. & \ref{sec:track-b}\\
Learnable hybrid basis & Can the network learn its own B-spline/RBF blend? & The gate works as designed but never beats the better pure basis, and it overfits on the pendulum past 3000 epochs. & \ref{sec:track-c}\\
Stiffness--solver map & Do explicit solvers fail as the pendulum's damping grows? & Zero divergences, as linear stability predicts. The failures trace to the random seed and a spurious learned equilibrium. & \ref{sec:track-d}\\
Sparse regression and epidemic transfer & How does KAN-ODE compare with SINDy under noise, and do the SIR fixes transfer to a non-mechanistic outbreak curve? & KAN-ODE is more accurate at every noise level, but degrades $1540\times$ against SINDy's $9\times$. Pruning finds no sparse structure. Two of three SIR fixes transfer. & \ref{sec:track-e}\\
Epoch-budget validation & Do the $10^4$-epoch conclusions survive five times the budget? & Solver and basis rankings hold. KAN vs.\ MLP \emph{reverses}: the MLP is $6.8\times$ better at $5\times10^4$ epochs. & \ref{sec:phase4}\\
Comparison with the base paper & Audit our recipe against the paper field by field. & Architecture and basis match. The KAN's training speed is reproduced within $1.4\times$. The tenfold speed advantage depends on the MLP baseline's optimizer settings. & \ref{sec:audit}\\
\bottomrule
\end{tabularx}
\end{table*}

The rest of the paper follows the course's required structure.
Section~\ref{sec:method} gives the numerical methods: the KAN-ODE model, the
integrators, the basis families, the two gradient methods, the hybrid basis
and sparse regression. Section~\ref{sec:impl} describes the datasets,
training recipe, evaluation protocol and compute budget.
Section~\ref{sec:results} reports each study in the order of
Table~\ref{tab:work-summary}, compares our findings with the base paper, and
discusses the patterns that recur across studies.
Section~\ref{sec:conclusion} concludes with limitations and future work.
